\documentclass[11pt]{article}
\usepackage[margin=0.9in]{geometry}
\usepackage{booktabs,amsmath,float}
\usepackage[T1]{fontenc}\usepackage{lmodern}
\usepackage[colorlinks=true,linkcolor=black,urlcolor=black]{hyperref}
\title{Note B: baseline and robustness results\\ \large Unemployment and prices after a monetary policy surprise}
\author{Working tables (generated by \texttt{baseline/make\_report.py})}
\date{}
\begin{document}\maketitle

\section*{What the tables show}
Each table reports the response of the unemployment gap (percentage points) and of the core PCE price \emph{level} (percent) to a monetary policy surprise that raises the 2-year Treasury yield by 1 percentage point in the announcement window, at 12, 24 and 36 months. Responses are traced through a daily block (the 2-year yield, and the market financial conditions index (FCI) of Note A where included) into a monthly system for the unemployment gap and core PCE inflation, with 3, 6 or 12 monthly lags. In every row the 2-year yield is shown without and with the FCI, and the FCI is in the same form as the rate: nominal, real or r*-gap.

\paragraph{Common specification.} Surprises: the Nakamura--Steinsson policy news series, 252 FOMC announcements from January 1996 to July 2026, identified in the announcement window with macro variables held fixed; macro variables first within the month; no deterministic trend in any equation; sample to September 2026 unless stated. Intervals are 90\% bootstrap intervals from 200 draws (residual resampling in the monthly equations and resampling of announcements; the daily block is not re-estimated). A star marks a 90\% interval that excludes zero.

\paragraph{Cases.} \textbf{Baseline}: start February 1994, all data in. \textbf{Robustness 1} (sample; r*-gap only): start 1973; March--December 2020 dropped from all equations, start 1994 and 1973; sample ending December 2019, start 1994 and 1973. \textbf{Robustness 2} (r* adjustment; start 1994 and 1973): Kiley's r*, and the 2-year yield minus its trailing 8-year average. \textbf{Robustness 3} (r*-gap; all data): the Baa--10-year spread added to the daily block, and the excess bond premium (EBP) added as a third monthly variable.

\paragraph{Choices made in building these tables} (to be confirmed).
\begin{itemize}\itemsep0pt
\item ``Dropping 2020'' removes the monthly observations and the daily observations for March through December 2020 from the estimation of every equation. The 2020 values still enter as lagged regressors in the following months, and the 2020 announcements remain among the shocks.
\item Robustness 1 has five cases (one start-date change, two drop-2020, two end-2019).
\item In Robustness 2 the FCI is held at Note A's r*-gap index (Laubach--Williams) for both the Kiley and the 8-year-average adjustments. The 8-year average needs eight years of daily data, so its daily block is estimated from 1998.
\item The Baa-spread case exists only with the FCI. The EBP source file is undocumented [unverified].
\item Starts before 1990 use monthly splices for the 2-year yield, expected inflation and the FCI before the daily data begin (FS-A29 in \texttt{FS\_QA.md}).
\end{itemize}

\section{Baseline}
\begin{table}[p]\centering\footnotesize\setlength{\tabcolsep}{2.2pt}
\caption{Baseline: 1994 to the present, all data in, no trend}\label{tab:B}
\begin{tabular}{llcrrrrrrr}\toprule
 & & & & \multicolumn{3}{c}{Unemployment gap (pp)} & \multicolumn{3}{c}{Price level (\%)}\\ \cmidrule(lr){5-7}\cmidrule(lr){8-10}
2-year & FCI & Lags & Obs & 12 m & 24 m & 36 m & 12 m & 24 m & 36 m\\ \midrule
Nominal & none & 3 & 392 & $-$0.53$^{*}$ & $-$0.37$^{*}$ & $-$0.26$^{*}$ & 0.10 & 0.07 & 0.04\\
 & & & & {\scriptsize [$-$0.72, $-$0.27]} & {\scriptsize [$-$0.62, $-$0.07]} & {\scriptsize [$-$0.53, $-$0.01]} & {\scriptsize [$-$0.07, 0.28]} & {\scriptsize [$-$0.20, 0.35]} & {\scriptsize [$-$0.30, 0.40]}\\
 &  & 6 & 392 & $-$0.48$^{*}$ & $-$0.34$^{*}$ & $-$0.22 & 0.09 & 0.03 & $-$0.08\\
 & & & & {\scriptsize [$-$0.69, $-$0.20]} & {\scriptsize [$-$0.60, $-$0.05]} & {\scriptsize [$-$0.47, 0.13]} & {\scriptsize [$-$0.11, 0.28]} & {\scriptsize [$-$0.33, 0.36]} & {\scriptsize [$-$0.55, 0.37]}\\
 &  & 12 & 392 & $-$0.70$^{*}$ & $-$0.38$^{*}$ & $-$0.20 & 0.08 & 0.01 & $-$0.12\\
 & & & & {\scriptsize [$-$1.17, $-$0.28]} & {\scriptsize [$-$0.72, $-$0.08]} & {\scriptsize [$-$0.56, 0.09]} & {\scriptsize [$-$0.10, 0.26]} & {\scriptsize [$-$0.30, 0.34]} & {\scriptsize [$-$0.55, 0.38]}\\
Nominal & nominal & 3 & 392 & 0.03 & 0.18 & 0.18 & $-$0.06 & $-$0.21 & $-$0.31\\
 & & & & {\scriptsize [$-$0.21, 0.28]} & {\scriptsize [$-$0.09, 0.38]} & {\scriptsize [0.00, 0.31]} & {\scriptsize [$-$0.24, 0.14]} & {\scriptsize [$-$0.48, 0.12]} & {\scriptsize [$-$0.67, 0.15]}\\
 &  & 6 & 392 & 0.03 & 0.24 & 0.27$^{*}$ & $-$0.08 & $-$0.29 & $-$0.48\\
 & & & & {\scriptsize [$-$0.24, 0.21]} & {\scriptsize [$-$0.03, 0.40]} & {\scriptsize [0.05, 0.44]} & {\scriptsize [$-$0.30, 0.11]} & {\scriptsize [$-$0.61, 0.06]} & {\scriptsize [$-$1.01, 0.06]}\\
 &  & 12 & 392 & $-$0.02 & 0.24 & 0.32$^{*}$ & $-$0.09 & $-$0.33 & $-$0.57$^{*}$\\
 & & & & {\scriptsize [$-$0.37, 0.39]} & {\scriptsize [$-$0.04, 0.53]} & {\scriptsize [0.07, 0.55]} & {\scriptsize [$-$0.32, 0.12]} & {\scriptsize [$-$0.77, 0.03]} & {\scriptsize [$-$1.19, $-$0.09]}\\
\addlinespace
Real & none & 3 & 392 & $-$0.32$^{*}$ & $-$0.30$^{*}$ & $-$0.25 & $-$0.03 & $-$0.05 & $-$0.06\\
 & & & & {\scriptsize [$-$0.53, $-$0.09]} & {\scriptsize [$-$0.55, $-$0.01]} & {\scriptsize [$-$0.49, 0.03]} & {\scriptsize [$-$0.18, 0.10]} & {\scriptsize [$-$0.28, 0.16]} & {\scriptsize [$-$0.36, 0.21]}\\
 &  & 6 & 392 & $-$0.31$^{*}$ & $-$0.26$^{*}$ & $-$0.19 & 0.00 & $-$0.06 & $-$0.16\\
 & & & & {\scriptsize [$-$0.54, $-$0.11]} & {\scriptsize [$-$0.52, $-$0.02]} & {\scriptsize [$-$0.46, 0.07]} & {\scriptsize [$-$0.16, 0.18]} & {\scriptsize [$-$0.30, 0.27]} & {\scriptsize [$-$0.52, 0.27]}\\
 &  & 12 & 392 & $-$0.54$^{*}$ & $-$0.28 & $-$0.16 & $-$0.03 & $-$0.10 & $-$0.21\\
 & & & & {\scriptsize [$-$0.92, $-$0.20]} & {\scriptsize [$-$0.58, 0.00]} & {\scriptsize [$-$0.50, 0.13]} & {\scriptsize [$-$0.20, 0.13]} & {\scriptsize [$-$0.40, 0.19]} & {\scriptsize [$-$0.63, 0.26]}\\
Real & real & 3 & 392 & 0.11 & 0.11 & 0.08 & $-$0.14 & $-$0.24 & $-$0.30\\
 & & & & {\scriptsize [$-$0.18, 0.40]} & {\scriptsize [$-$0.17, 0.35]} & {\scriptsize [$-$0.12, 0.22]} & {\scriptsize [$-$0.33, 0.06]} & {\scriptsize [$-$0.53, 0.08]} & {\scriptsize [$-$0.67, 0.11]}\\
 &  & 6 & 392 & 0.11 & 0.22 & 0.20$^{*}$ & $-$0.14 & $-$0.33 & $-$0.50\\
 & & & & {\scriptsize [$-$0.08, 0.38]} & {\scriptsize [$-$0.01, 0.54]} & {\scriptsize [0.02, 0.46]} & {\scriptsize [$-$0.31, 0.06]} & {\scriptsize [$-$0.63, 0.03]} & {\scriptsize [$-$0.92, 0.07]}\\
 &  & 12 & 392 & 0.07 & 0.25 & 0.26 & $-$0.18$^{*}$ & $-$0.38 & $-$0.54$^{*}$\\
 & & & & {\scriptsize [$-$0.30, 0.43]} & {\scriptsize [$-$0.05, 0.48]} & {\scriptsize [$-$0.02, 0.47]} & {\scriptsize [$-$0.38, 0.00]} & {\scriptsize [$-$0.74, 0.03]} & {\scriptsize [$-$1.14, 0.00]}\\
\addlinespace
r*-gap & none & 3 & 392 & $-$0.17 & $-$0.16 & $-$0.13 & $-$0.09 & $-$0.11 & $-$0.13\\
 & & & & {\scriptsize [$-$0.48, 0.06]} & {\scriptsize [$-$0.53, 0.13]} & {\scriptsize [$-$0.51, 0.15]} & {\scriptsize [$-$0.25, 0.07]} & {\scriptsize [$-$0.39, 0.17]} & {\scriptsize [$-$0.49, 0.26]}\\
 &  & 6 & 392 & $-$0.11 & $-$0.08 & $-$0.03 & $-$0.12 & $-$0.28 & $-$0.44\\
 & & & & {\scriptsize [$-$0.38, 0.13]} & {\scriptsize [$-$0.36, 0.20]} & {\scriptsize [$-$0.38, 0.26]} & {\scriptsize [$-$0.28, 0.06]} & {\scriptsize [$-$0.62, 0.08]} & {\scriptsize [$-$0.95, 0.09]}\\
 &  & 12 & 392 & $-$0.28 & $-$0.06 & 0.04 & $-$0.18 & $-$0.39$^{*}$ & $-$0.62$^{*}$\\
 & & & & {\scriptsize [$-$0.77, 0.08]} & {\scriptsize [$-$0.54, 0.30]} & {\scriptsize [$-$0.41, 0.45]} & {\scriptsize [$-$0.34, 0.03]} & {\scriptsize [$-$0.68, 0.00]} & {\scriptsize [$-$1.11, $-$0.06]}\\
r*-gap & r*-gap & 3 & 392 & 0.25 & 0.17 & 0.06 & $-$0.17 & $-$0.21 & $-$0.22\\
 & & & & {\scriptsize [$-$0.03, 0.50]} & {\scriptsize [$-$0.12, 0.43]} & {\scriptsize [$-$0.16, 0.24]} & {\scriptsize [$-$0.40, 0.07]} & {\scriptsize [$-$0.57, 0.23]} & {\scriptsize [$-$0.67, 0.36]}\\
 &  & 6 & 392 & 0.31$^{*}$ & 0.34$^{*}$ & 0.25$^{*}$ & $-$0.27$^{*}$ & $-$0.53$^{*}$ & $-$0.73$^{*}$\\
 & & & & {\scriptsize [0.07, 0.63]} & {\scriptsize [0.10, 0.68]} & {\scriptsize [0.04, 0.53]} & {\scriptsize [$-$0.48, $-$0.06]} & {\scriptsize [$-$0.97, $-$0.08]} & {\scriptsize [$-$1.42, $-$0.09]}\\
 &  & 12 & 392 & 0.31 & 0.47$^{*}$ & 0.41$^{*}$ & $-$0.36$^{*}$ & $-$0.71$^{*}$ & $-$0.97$^{*}$\\
 & & & & {\scriptsize [$-$0.03, 0.72]} & {\scriptsize [0.15, 0.75]} & {\scriptsize [0.09, 0.70]} & {\scriptsize [$-$0.54, $-$0.20]} & {\scriptsize [$-$1.11, $-$0.34]} & {\scriptsize [$-$1.64, $-$0.33]}\\
\bottomrule\end{tabular}
\par\smallskip\parbox{0.96\linewidth}{\scriptsize Nominal, real and r*-gap 2-year yield; the FCI is the matching index of Note A in each case. Per +1 pp on the 2-year yield. $^{*}$ = 90\% interval excludes zero; intervals in brackets. Obs = monthly observations in the equations.}
\end{table}


\section{Robustness 1: sample (r*-gap only)}
\begin{table}[p]\centering\footnotesize\setlength{\tabcolsep}{2.2pt}
\caption{Robustness 1a: monthly equations start in 1973}\label{tab:R1a}
\begin{tabular}{llcrrrrrrr}\toprule
 & & & & \multicolumn{3}{c}{Unemployment gap (pp)} & \multicolumn{3}{c}{Price level (\%)}\\ \cmidrule(lr){5-7}\cmidrule(lr){8-10}
2-year & FCI & Lags & Obs & 12 m & 24 m & 36 m & 12 m & 24 m & 36 m\\ \midrule
r*-gap & none & 3 & 642 & 0.01 & 0.05 & 0.05 & $-$0.09 & $-$0.15 & $-$0.19\\
 & & & & {\scriptsize [$-$0.10, 0.20]} & {\scriptsize [$-$0.10, 0.26]} & {\scriptsize [$-$0.10, 0.27]} & {\scriptsize [$-$0.23, 0.06]} & {\scriptsize [$-$0.48, 0.22]} & {\scriptsize [$-$0.71, 0.38]}\\
 &  & 6 & 639 & 0.00 & 0.03 & 0.03 & $-$0.12 & $-$0.31 & $-$0.52\\
 & & & & {\scriptsize [$-$0.13, 0.12]} & {\scriptsize [$-$0.14, 0.18]} & {\scriptsize [$-$0.15, 0.17]} & {\scriptsize [$-$0.24, 0.02]} & {\scriptsize [$-$0.60, 0.02]} & {\scriptsize [$-$1.02, 0.01]}\\
 &  & 12 & 633 & $-$0.05 & 0.00 & 0.00 & $-$0.23$^{*}$ & $-$0.50$^{*}$ & $-$0.78$^{*}$\\
 & & & & {\scriptsize [$-$0.22, 0.13]} & {\scriptsize [$-$0.19, 0.17]} & {\scriptsize [$-$0.19, 0.24]} & {\scriptsize [$-$0.36, $-$0.11]} & {\scriptsize [$-$0.79, $-$0.24]} & {\scriptsize [$-$1.25, $-$0.31]}\\
r*-gap & r*-gap & 3 & 642 & 0.07 & 0.03 & 0.00 & $-$0.08 & $-$0.15 & $-$0.23\\
 & & & & {\scriptsize [$-$0.05, 0.18]} & {\scriptsize [$-$0.13, 0.17]} & {\scriptsize [$-$0.16, 0.16]} & {\scriptsize [$-$0.23, 0.07]} & {\scriptsize [$-$0.51, 0.18]} & {\scriptsize [$-$0.75, 0.28]}\\
 &  & 6 & 639 & 0.10 & 0.08 & 0.07 & $-$0.14 & $-$0.35$^{*}$ & $-$0.59$^{*}$\\
 & & & & {\scriptsize [$-$0.06, 0.28]} & {\scriptsize [$-$0.11, 0.28]} & {\scriptsize [$-$0.13, 0.27]} & {\scriptsize [$-$0.29, 0.02]} & {\scriptsize [$-$0.67, $-$0.01]} & {\scriptsize [$-$1.11, $-$0.14]}\\
 &  & 12 & 633 & 0.18 & 0.15 & 0.14 & $-$0.27$^{*}$ & $-$0.62$^{*}$ & $-$0.94$^{*}$\\
 & & & & {\scriptsize [$-$0.04, 0.37]} & {\scriptsize [$-$0.01, 0.31]} & {\scriptsize [$-$0.04, 0.32]} & {\scriptsize [$-$0.44, $-$0.12]} & {\scriptsize [$-$0.94, $-$0.29]} & {\scriptsize [$-$1.46, $-$0.42]}\\
\bottomrule\end{tabular}
\par\smallskip\parbox{0.96\linewidth}{\scriptsize r*-gap only. Pre-1990 months use the spliced monthly series (FS-A29); shocks still 1996 on. Per +1 pp on the 2-year yield. $^{*}$ = 90\% interval excludes zero; intervals in brackets. Obs = monthly observations in the equations.}
\end{table}
\begin{table}[p]\centering\footnotesize\setlength{\tabcolsep}{2.2pt}
\caption{Robustness 1b: 2020 dropped from all equations, start 1994}\label{tab:R1b}
\begin{tabular}{llcrrrrrrr}\toprule
 & & & & \multicolumn{3}{c}{Unemployment gap (pp)} & \multicolumn{3}{c}{Price level (\%)}\\ \cmidrule(lr){5-7}\cmidrule(lr){8-10}
2-year & FCI & Lags & Obs & 12 m & 24 m & 36 m & 12 m & 24 m & 36 m\\ \midrule
r*-gap & none & 3 & 382 & 0.08 & 0.16 & 0.20 & $-$0.15 & $-$0.22 & $-$0.28\\
 & & & & {\scriptsize [$-$0.08, 0.26]} & {\scriptsize [$-$0.09, 0.42]} & {\scriptsize [$-$0.10, 0.50]} & {\scriptsize [$-$0.33, 0.03]} & {\scriptsize [$-$0.59, 0.15]} & {\scriptsize [$-$0.79, 0.23]}\\
 &  & 6 & 382 & 0.03 & 0.04 & 0.05 & $-$0.12 & $-$0.26 & $-$0.39\\
 & & & & {\scriptsize [$-$0.17, 0.19]} & {\scriptsize [$-$0.28, 0.31]} & {\scriptsize [$-$0.36, 0.39]} & {\scriptsize [$-$0.30, 0.04]} & {\scriptsize [$-$0.56, 0.11]} & {\scriptsize [$-$0.88, 0.19]}\\
 &  & 12 & 382 & $-$0.08 & $-$0.15 & $-$0.16 & $-$0.18$^{*}$ & $-$0.31 & $-$0.44\\
 & & & & {\scriptsize [$-$0.26, 0.14]} & {\scriptsize [$-$0.46, 0.22]} & {\scriptsize [$-$0.57, 0.31]} & {\scriptsize [$-$0.35, $-$0.02]} & {\scriptsize [$-$0.66, 0.07]} & {\scriptsize [$-$0.97, 0.12]}\\
r*-gap & r*-gap & 3 & 382 & 0.30$^{*}$ & 0.38$^{*}$ & 0.34$^{*}$ & $-$0.18 & $-$0.28 & $-$0.37\\
 & & & & {\scriptsize [0.13, 0.47]} & {\scriptsize [0.13, 0.61]} & {\scriptsize [0.05, 0.59]} & {\scriptsize [$-$0.35, 0.02]} & {\scriptsize [$-$0.67, 0.08]} & {\scriptsize [$-$0.82, 0.19]}\\
 &  & 6 & 382 & 0.32$^{*}$ & 0.39$^{*}$ & 0.38$^{*}$ & $-$0.25$^{*}$ & $-$0.48$^{*}$ & $-$0.68$^{*}$\\
 & & & & {\scriptsize [0.15, 0.49]} & {\scriptsize [0.11, 0.70]} & {\scriptsize [0.05, 0.77]} & {\scriptsize [$-$0.45, $-$0.04]} & {\scriptsize [$-$0.92, $-$0.07]} & {\scriptsize [$-$1.36, $-$0.07]}\\
 &  & 12 & 382 & 0.35$^{*}$ & 0.41$^{*}$ & 0.37$^{*}$ & $-$0.33$^{*}$ & $-$0.61$^{*}$ & $-$0.84$^{*}$\\
 & & & & {\scriptsize [0.15, 0.55]} & {\scriptsize [0.15, 0.70]} & {\scriptsize [0.03, 0.73]} & {\scriptsize [$-$0.51, $-$0.12]} & {\scriptsize [$-$1.02, $-$0.15]} & {\scriptsize [$-$1.51, $-$0.15]}\\
\bottomrule\end{tabular}
\par\smallskip\parbox{0.96\linewidth}{\scriptsize r*-gap only. March to December 2020 dropped from the monthly and the daily equations. Per +1 pp on the 2-year yield. $^{*}$ = 90\% interval excludes zero; intervals in brackets. Obs = monthly observations in the equations.}
\end{table}
\begin{table}[p]\centering\footnotesize\setlength{\tabcolsep}{2.2pt}
\caption{Robustness 1c: 2020 dropped from all equations, start 1973}\label{tab:R1c}
\begin{tabular}{llcrrrrrrr}\toprule
 & & & & \multicolumn{3}{c}{Unemployment gap (pp)} & \multicolumn{3}{c}{Price level (\%)}\\ \cmidrule(lr){5-7}\cmidrule(lr){8-10}
2-year & FCI & Lags & Obs & 12 m & 24 m & 36 m & 12 m & 24 m & 36 m\\ \midrule
r*-gap & none & 3 & 632 & 0.06 & 0.11 & 0.13 & $-$0.11 & $-$0.20 & $-$0.27\\
 & & & & {\scriptsize [$-$0.02, 0.14]} & {\scriptsize [$-$0.04, 0.24]} & {\scriptsize [$-$0.04, 0.30]} & {\scriptsize [$-$0.26, 0.03]} & {\scriptsize [$-$0.53, 0.13]} & {\scriptsize [$-$0.80, 0.22]}\\
 &  & 6 & 629 & 0.11$^{*}$ & 0.10 & 0.04 & $-$0.17$^{*}$ & $-$0.39$^{*}$ & $-$0.62$^{*}$\\
 & & & & {\scriptsize [0.01, 0.22]} & {\scriptsize [$-$0.06, 0.24]} & {\scriptsize [$-$0.15, 0.22]} & {\scriptsize [$-$0.30, $-$0.03]} & {\scriptsize [$-$0.70, $-$0.05]} & {\scriptsize [$-$1.14, $-$0.12]}\\
 &  & 12 & 623 & 0.18$^{*}$ & 0.14 & 0.04 & $-$0.29$^{*}$ & $-$0.64$^{*}$ & $-$0.96$^{*}$\\
 & & & & {\scriptsize [0.07, 0.29]} & {\scriptsize [0.00, 0.30]} & {\scriptsize [$-$0.14, 0.22]} & {\scriptsize [$-$0.41, $-$0.16]} & {\scriptsize [$-$0.89, $-$0.37]} & {\scriptsize [$-$1.35, $-$0.58]}\\
r*-gap & r*-gap & 3 & 632 & 0.12$^{*}$ & 0.10 & 0.08 & $-$0.09 & $-$0.24$^{*}$ & $-$0.44$^{*}$\\
 & & & & {\scriptsize [0.03, 0.22]} & {\scriptsize [$-$0.04, 0.25]} & {\scriptsize [$-$0.08, 0.24]} & {\scriptsize [$-$0.25, 0.00]} & {\scriptsize [$-$0.57, $-$0.01]} & {\scriptsize [$-$0.90, $-$0.05]}\\
 &  & 6 & 629 & 0.27$^{*}$ & 0.16$^{*}$ & 0.07 & $-$0.20$^{*}$ & $-$0.48$^{*}$ & $-$0.77$^{*}$\\
 & & & & {\scriptsize [0.16, 0.39]} & {\scriptsize [0.00, 0.32]} & {\scriptsize [$-$0.12, 0.26]} & {\scriptsize [$-$0.39, $-$0.07]} & {\scriptsize [$-$0.90, $-$0.18]} & {\scriptsize [$-$1.37, $-$0.30]}\\
 &  & 12 & 623 & 0.32$^{*}$ & 0.24$^{*}$ & 0.16 & $-$0.32$^{*}$ & $-$0.74$^{*}$ & $-$1.11$^{*}$\\
 & & & & {\scriptsize [0.20, 0.43]} & {\scriptsize [0.10, 0.39]} & {\scriptsize [0.00, 0.32]} & {\scriptsize [$-$0.46, $-$0.16]} & {\scriptsize [$-$1.05, $-$0.41]} & {\scriptsize [$-$1.69, $-$0.56]}\\
\bottomrule\end{tabular}
\par\smallskip\parbox{0.96\linewidth}{\scriptsize r*-gap only. Per +1 pp on the 2-year yield. $^{*}$ = 90\% interval excludes zero; intervals in brackets. Obs = monthly observations in the equations.}
\end{table}
\begin{table}[p]\centering\footnotesize\setlength{\tabcolsep}{2.2pt}
\caption{Robustness 1d: sample ends in December 2019, start 1994}\label{tab:R1d}
\begin{tabular}{llcrrrrrrr}\toprule
 & & & & \multicolumn{3}{c}{Unemployment gap (pp)} & \multicolumn{3}{c}{Price level (\%)}\\ \cmidrule(lr){5-7}\cmidrule(lr){8-10}
2-year & FCI & Lags & Obs & 12 m & 24 m & 36 m & 12 m & 24 m & 36 m\\ \midrule
r*-gap & none & 3 & 311 & $-$0.09 & $-$0.13 & $-$0.15 & 0.02 & 0.08 & 0.13\\
 & & & & {\scriptsize [$-$0.26, 0.12]} & {\scriptsize [$-$0.43, 0.18]} & {\scriptsize [$-$0.50, 0.25]} & {\scriptsize [$-$0.05, 0.11]} & {\scriptsize [$-$0.04, 0.22]} & {\scriptsize [$-$0.04, 0.30]}\\
 &  & 6 & 311 & $-$0.09 & $-$0.14 & $-$0.15 & 0.04 & 0.08 & 0.10\\
 & & & & {\scriptsize [$-$0.33, 0.13]} & {\scriptsize [$-$0.51, 0.23]} & {\scriptsize [$-$0.59, 0.36]} & {\scriptsize [$-$0.05, 0.14]} & {\scriptsize [$-$0.12, 0.24]} & {\scriptsize [$-$0.19, 0.33]}\\
 &  & 12 & 311 & $-$0.14 & $-$0.19 & $-$0.16 & 0.10 & 0.10 & 0.07\\
 & & & & {\scriptsize [$-$0.34, 0.08]} & {\scriptsize [$-$0.55, 0.17]} & {\scriptsize [$-$0.66, 0.29]} & {\scriptsize [$-$0.02, 0.22]} & {\scriptsize [$-$0.07, 0.28]} & {\scriptsize [$-$0.19, 0.28]}\\
r*-gap & r*-gap & 3 & 311 & 0.20$^{*}$ & 0.21 & 0.17 & $-$0.04 & 0.00 & 0.03\\
 & & & & {\scriptsize [0.00, 0.44]} & {\scriptsize [$-$0.10, 0.54]} & {\scriptsize [$-$0.17, 0.57]} & {\scriptsize [$-$0.14, 0.07]} & {\scriptsize [$-$0.17, 0.20]} & {\scriptsize [$-$0.19, 0.28]}\\
 &  & 6 & 311 & 0.23$^{*}$ & 0.21 & 0.17 & $-$0.07 & $-$0.04 & $-$0.04\\
 & & & & {\scriptsize [0.01, 0.48]} & {\scriptsize [$-$0.14, 0.66]} & {\scriptsize [$-$0.27, 0.74]} & {\scriptsize [$-$0.19, 0.05]} & {\scriptsize [$-$0.23, 0.17]} & {\scriptsize [$-$0.26, 0.23]}\\
 &  & 12 & 311 & 0.30$^{*}$ & 0.30 & 0.23 & $-$0.05 & $-$0.12 & $-$0.21\\
 & & & & {\scriptsize [0.12, 0.49]} & {\scriptsize [$-$0.04, 0.59]} & {\scriptsize [$-$0.21, 0.61]} & {\scriptsize [$-$0.24, 0.07]} & {\scriptsize [$-$0.42, 0.07]} & {\scriptsize [$-$0.65, 0.07]}\\
\bottomrule\end{tabular}
\par\smallskip\parbox{0.96\linewidth}{\scriptsize r*-gap only. All equations, monthly and daily. Per +1 pp on the 2-year yield. $^{*}$ = 90\% interval excludes zero; intervals in brackets. Obs = monthly observations in the equations.}
\end{table}
\begin{table}[p]\centering\footnotesize\setlength{\tabcolsep}{2.2pt}
\caption{Robustness 1e: sample ends in December 2019, start 1973}\label{tab:R1e}
\begin{tabular}{llcrrrrrrr}\toprule
 & & & & \multicolumn{3}{c}{Unemployment gap (pp)} & \multicolumn{3}{c}{Price level (\%)}\\ \cmidrule(lr){5-7}\cmidrule(lr){8-10}
2-year & FCI & Lags & Obs & 12 m & 24 m & 36 m & 12 m & 24 m & 36 m\\ \midrule
r*-gap & none & 3 & 561 & 0.03 & 0.05 & 0.06 & $-$0.05 & $-$0.06 & $-$0.07\\
 & & & & {\scriptsize [$-$0.06, 0.13]} & {\scriptsize [$-$0.09, 0.21]} & {\scriptsize [$-$0.10, 0.24]} & {\scriptsize [$-$0.19, 0.12]} & {\scriptsize [$-$0.45, 0.37]} & {\scriptsize [$-$0.67, 0.60]}\\
 &  & 6 & 558 & 0.11$^{*}$ & 0.09 & 0.04 & $-$0.10 & $-$0.27 & $-$0.46\\
 & & & & {\scriptsize [0.01, 0.24]} & {\scriptsize [$-$0.05, 0.28]} & {\scriptsize [$-$0.12, 0.23]} & {\scriptsize [$-$0.24, 0.05]} & {\scriptsize [$-$0.58, 0.07]} & {\scriptsize [$-$1.06, 0.10]}\\
 &  & 12 & 552 & 0.20$^{*}$ & 0.18$^{*}$ & 0.07 & $-$0.22$^{*}$ & $-$0.51$^{*}$ & $-$0.83$^{*}$\\
 & & & & {\scriptsize [0.07, 0.35]} & {\scriptsize [0.01, 0.39]} & {\scriptsize [$-$0.14, 0.35]} & {\scriptsize [$-$0.35, $-$0.07]} & {\scriptsize [$-$0.82, $-$0.16]} & {\scriptsize [$-$1.36, $-$0.25]}\\
r*-gap & r*-gap & 3 & 561 & 0.08$^{*}$ & 0.04 & 0.00 & $-$0.01 & $-$0.06 & $-$0.19\\
 & & & & {\scriptsize [0.00, 0.18]} & {\scriptsize [$-$0.11, 0.19]} & {\scriptsize [$-$0.18, 0.19]} & {\scriptsize [$-$0.14, 0.14]} & {\scriptsize [$-$0.33, 0.24]} & {\scriptsize [$-$0.60, 0.25]}\\
 &  & 6 & 558 & 0.26$^{*}$ & 0.13 & 0.02 & $-$0.13 & $-$0.34 & $-$0.59$^{*}$\\
 & & & & {\scriptsize [0.15, 0.39]} & {\scriptsize [$-$0.04, 0.28]} & {\scriptsize [$-$0.15, 0.20]} & {\scriptsize [$-$0.28, 0.01]} & {\scriptsize [$-$0.67, 0.02]} & {\scriptsize [$-$1.12, $-$0.06]}\\
 &  & 12 & 552 & 0.33$^{*}$ & 0.25$^{*}$ & 0.15 & $-$0.23$^{*}$ & $-$0.54$^{*}$ & $-$0.90$^{*}$\\
 & & & & {\scriptsize [0.21, 0.47]} & {\scriptsize [0.10, 0.44]} & {\scriptsize [$-$0.03, 0.36]} & {\scriptsize [$-$0.41, $-$0.07]} & {\scriptsize [$-$0.94, $-$0.21]} & {\scriptsize [$-$1.53, $-$0.41]}\\
\bottomrule\end{tabular}
\par\smallskip\parbox{0.96\linewidth}{\scriptsize r*-gap only. Per +1 pp on the 2-year yield. $^{*}$ = 90\% interval excludes zero; intervals in brackets. Obs = monthly observations in the equations.}
\end{table}


\section{Robustness 2: the r* adjustment}
\begin{table}[p]\centering\footnotesize\setlength{\tabcolsep}{2.2pt}
\caption{Robustness 2a: Kiley r*, start 1994}\label{tab:R2a}
\begin{tabular}{llcrrrrrrr}\toprule
 & & & & \multicolumn{3}{c}{Unemployment gap (pp)} & \multicolumn{3}{c}{Price level (\%)}\\ \cmidrule(lr){5-7}\cmidrule(lr){8-10}
2-year & FCI & Lags & Obs & 12 m & 24 m & 36 m & 12 m & 24 m & 36 m\\ \midrule
r*-gap (Kiley r*) & none & 3 & 392 & $-$0.31$^{*}$ & $-$0.32$^{*}$ & $-$0.27 & $-$0.05 & $-$0.03 & $-$0.02\\
 & & & & {\scriptsize [$-$0.50, $-$0.05]} & {\scriptsize [$-$0.55, $-$0.03]} & {\scriptsize [$-$0.50, 0.00]} & {\scriptsize [$-$0.23, 0.10]} & {\scriptsize [$-$0.28, 0.21]} & {\scriptsize [$-$0.34, 0.32]}\\
 &  & 6 & 392 & $-$0.27$^{*}$ & $-$0.26$^{*}$ & $-$0.20 & $-$0.05 & $-$0.10 & $-$0.17\\
 & & & & {\scriptsize [$-$0.64, $-$0.06]} & {\scriptsize [$-$0.58, 0.00]} & {\scriptsize [$-$0.54, 0.07]} & {\scriptsize [$-$0.21, 0.14]} & {\scriptsize [$-$0.41, 0.22]} & {\scriptsize [$-$0.64, 0.29]}\\
 &  & 12 & 392 & $-$0.47 & $-$0.26 & $-$0.17 & $-$0.09 & $-$0.11 & $-$0.14\\
 & & & & {\scriptsize [$-$0.91, 0.01]} & {\scriptsize [$-$0.59, 0.17]} & {\scriptsize [$-$0.54, 0.26]} & {\scriptsize [$-$0.26, 0.12]} & {\scriptsize [$-$0.49, 0.24]} & {\scriptsize [$-$0.67, 0.35]}\\
r*-gap (Kiley r*) & r*-gap & 3 & 392 & 0.15 & 0.07 & 0.01 & $-$0.13 & $-$0.15 & $-$0.15\\
 & & & & {\scriptsize [$-$0.11, 0.43]} & {\scriptsize [$-$0.18, 0.32]} & {\scriptsize [$-$0.19, 0.17]} & {\scriptsize [$-$0.34, 0.06]} & {\scriptsize [$-$0.50, 0.18]} & {\scriptsize [$-$0.58, 0.27]}\\
 &  & 6 & 392 & 0.24 & 0.25 & 0.17 & $-$0.23$^{*}$ & $-$0.42 & $-$0.55\\
 & & & & {\scriptsize [$-$0.05, 0.55]} & {\scriptsize [$-$0.03, 0.54]} & {\scriptsize [$-$0.03, 0.38]} & {\scriptsize [$-$0.44, 0.00]} & {\scriptsize [$-$0.86, 0.04]} & {\scriptsize [$-$1.20, 0.08]}\\
 &  & 12 & 392 & 0.22 & 0.39$^{*}$ & 0.28$^{*}$ & $-$0.28$^{*}$ & $-$0.46$^{*}$ & $-$0.51\\
 & & & & {\scriptsize [$-$0.14, 0.61]} & {\scriptsize [0.08, 0.68]} & {\scriptsize [0.02, 0.56]} & {\scriptsize [$-$0.45, $-$0.11]} & {\scriptsize [$-$0.81, $-$0.12]} & {\scriptsize [$-$1.07, 0.02]}\\
\bottomrule\end{tabular}
\par\smallskip\parbox{0.96\linewidth}{\scriptsize 2-year minus expected inflation minus Kiley's r*; the FCI is the r*-gap index of Note A (Laubach-Williams). Per +1 pp on the 2-year yield. $^{*}$ = 90\% interval excludes zero; intervals in brackets. Obs = monthly observations in the equations.}
\end{table}
\begin{table}[p]\centering\footnotesize\setlength{\tabcolsep}{2.2pt}
\caption{Robustness 2b: Kiley r*, start 1973}\label{tab:R2b}
\begin{tabular}{llcrrrrrrr}\toprule
 & & & & \multicolumn{3}{c}{Unemployment gap (pp)} & \multicolumn{3}{c}{Price level (\%)}\\ \cmidrule(lr){5-7}\cmidrule(lr){8-10}
2-year & FCI & Lags & Obs & 12 m & 24 m & 36 m & 12 m & 24 m & 36 m\\ \midrule
r*-gap (Kiley r*) & none & 3 & 642 & $-$0.02 & 0.00 & 0.01 & $-$0.04 & $-$0.02 & 0.01\\
 & & & & {\scriptsize [$-$0.16, 0.09]} & {\scriptsize [$-$0.15, 0.18]} & {\scriptsize [$-$0.13, 0.18]} & {\scriptsize [$-$0.17, 0.10]} & {\scriptsize [$-$0.33, 0.29]} & {\scriptsize [$-$0.46, 0.52]}\\
 &  & 6 & 639 & $-$0.04 & $-$0.01 & 0.00 & $-$0.07 & $-$0.17 & $-$0.28\\
 & & & & {\scriptsize [$-$0.15, 0.07]} & {\scriptsize [$-$0.16, 0.16]} & {\scriptsize [$-$0.16, 0.15]} & {\scriptsize [$-$0.18, 0.04]} & {\scriptsize [$-$0.42, 0.08]} & {\scriptsize [$-$0.70, 0.15]}\\
 &  & 12 & 633 & $-$0.11 & $-$0.05 & $-$0.02 & $-$0.18$^{*}$ & $-$0.35$^{*}$ & $-$0.50$^{*}$\\
 & & & & {\scriptsize [$-$0.25, 0.03]} & {\scriptsize [$-$0.21, 0.12]} & {\scriptsize [$-$0.18, 0.14]} & {\scriptsize [$-$0.31, $-$0.06]} & {\scriptsize [$-$0.61, $-$0.11]} & {\scriptsize [$-$0.90, $-$0.10]}\\
r*-gap (Kiley r*) & r*-gap & 3 & 642 & 0.04 & $-$0.01 & $-$0.04 & $-$0.03 & 0.01 & 0.03\\
 & & & & {\scriptsize [$-$0.05, 0.15]} & {\scriptsize [$-$0.11, 0.10]} & {\scriptsize [$-$0.14, 0.06]} & {\scriptsize [$-$0.19, 0.11]} & {\scriptsize [$-$0.33, 0.34]} & {\scriptsize [$-$0.45, 0.47]}\\
 &  & 6 & 639 & 0.06 & 0.04 & 0.02 & $-$0.10 & $-$0.21 & $-$0.32\\
 & & & & {\scriptsize [$-$0.04, 0.19]} & {\scriptsize [$-$0.06, 0.16]} & {\scriptsize [$-$0.08, 0.14]} & {\scriptsize [$-$0.24, 0.04]} & {\scriptsize [$-$0.50, 0.06]} & {\scriptsize [$-$0.89, 0.09]}\\
 &  & 12 & 633 & 0.12 & 0.10 & 0.08 & $-$0.23$^{*}$ & $-$0.47$^{*}$ & $-$0.65$^{*}$\\
 & & & & {\scriptsize [$-$0.03, 0.28]} & {\scriptsize [$-$0.03, 0.24]} & {\scriptsize [$-$0.05, 0.21]} & {\scriptsize [$-$0.39, $-$0.09]} & {\scriptsize [$-$0.85, $-$0.17]} & {\scriptsize [$-$1.27, $-$0.18]}\\
\bottomrule\end{tabular}
\par\smallskip\parbox{0.96\linewidth}{\scriptsize As 2a. Per +1 pp on the 2-year yield. $^{*}$ = 90\% interval excludes zero; intervals in brackets. Obs = monthly observations in the equations.}
\end{table}
\begin{table}[p]\centering\footnotesize\setlength{\tabcolsep}{2.2pt}
\caption{Robustness 2c: 2-year minus its trailing 8-year average, start 1994}\label{tab:R2c}
\begin{tabular}{llcrrrrrrr}\toprule
 & & & & \multicolumn{3}{c}{Unemployment gap (pp)} & \multicolumn{3}{c}{Price level (\%)}\\ \cmidrule(lr){5-7}\cmidrule(lr){8-10}
2-year & FCI & Lags & Obs & 12 m & 24 m & 36 m & 12 m & 24 m & 36 m\\ \midrule
Minus 8-year average & none & 3 & 392 & $-$0.65$^{*}$ & $-$0.56$^{*}$ & $-$0.47$^{*}$ & 0.25$^{*}$ & 0.44$^{*}$ & 0.58$^{*}$\\
 & & & & {\scriptsize [$-$0.96, $-$0.27]} & {\scriptsize [$-$0.90, $-$0.14]} & {\scriptsize [$-$0.82, $-$0.09]} & {\scriptsize [0.05, 0.42]} & {\scriptsize [0.10, 0.80]} & {\scriptsize [0.11, 1.13]}\\
 &  & 6 & 392 & $-$0.56$^{*}$ & $-$0.48$^{*}$ & $-$0.39 & 0.13 & 0.20 & 0.23\\
 & & & & {\scriptsize [$-$0.93, $-$0.30]} & {\scriptsize [$-$0.94, $-$0.10]} & {\scriptsize [$-$0.86, 0.01]} & {\scriptsize [$-$0.06, 0.36]} & {\scriptsize [$-$0.24, 0.65]} & {\scriptsize [$-$0.42, 0.93]}\\
 &  & 12 & 392 & $-$0.74$^{*}$ & $-$0.48$^{*}$ & $-$0.34 & 0.11 & 0.14 & 0.14\\
 & & & & {\scriptsize [$-$1.22, $-$0.40]} & {\scriptsize [$-$0.92, $-$0.08]} & {\scriptsize [$-$0.85, 0.10]} & {\scriptsize [$-$0.12, 0.32]} & {\scriptsize [$-$0.28, 0.57]} & {\scriptsize [$-$0.54, 0.78]}\\
Minus 8-year average & r*-gap & 3 & 392 & $-$0.01 & 0.04 & 0.01 & 0.10 & 0.22 & 0.33\\
 & & & & {\scriptsize [$-$0.21, 0.22]} & {\scriptsize [$-$0.17, 0.25]} & {\scriptsize [$-$0.15, 0.15]} & {\scriptsize [$-$0.10, 0.29]} & {\scriptsize [$-$0.14, 0.61]} & {\scriptsize [$-$0.15, 0.87]}\\
 &  & 6 & 392 & 0.06 & 0.16 & 0.14 & $-$0.07 & $-$0.13 & $-$0.16\\
 & & & & {\scriptsize [$-$0.26, 0.33]} & {\scriptsize [$-$0.11, 0.41]} & {\scriptsize [$-$0.07, 0.35]} & {\scriptsize [$-$0.31, 0.14]} & {\scriptsize [$-$0.66, 0.34]} & {\scriptsize [$-$0.93, 0.53]}\\
 &  & 12 & 392 & 0.04 & 0.26$^{*}$ & 0.26$^{*}$ & $-$0.12 & $-$0.28 & $-$0.39\\
 & & & & {\scriptsize [$-$0.31, 0.45]} & {\scriptsize [0.07, 0.56]} & {\scriptsize [0.04, 0.53]} & {\scriptsize [$-$0.32, 0.05]} & {\scriptsize [$-$0.73, 0.13]} & {\scriptsize [$-$1.10, 0.37]}\\
\bottomrule\end{tabular}
\par\smallskip\parbox{0.96\linewidth}{\scriptsize The daily 8-year average needs eight years of daily data, so the daily block is estimated from 1998 on; the FCI is the r*-gap index of Note A. Per +1 pp on the 2-year yield. $^{*}$ = 90\% interval excludes zero; intervals in brackets. Obs = monthly observations in the equations.}
\end{table}
\begin{table}[p]\centering\footnotesize\setlength{\tabcolsep}{2.2pt}
\caption{Robustness 2d: real 2-year and real FCI with a quadratic trend, start 1973}\label{tab:R2d}
\begin{tabular}{llcrrrrrrr}\toprule
 & & & & \multicolumn{3}{c}{Unemployment gap (pp)} & \multicolumn{3}{c}{Price level (\%)}\\ \cmidrule(lr){5-7}\cmidrule(lr){8-10}
2-year & FCI & Lags & Obs & 12 m & 24 m & 36 m & 12 m & 24 m & 36 m\\ \midrule
Real & none & 3 & 642 & $-$0.06 & $-$0.03 & $-$0.02 & $-$0.04 & $-$0.07 & $-$0.10\\
 & & & & {\scriptsize [$-$0.14, 0.10]} & {\scriptsize [$-$0.15, 0.19]} & {\scriptsize [$-$0.14, 0.21]} & {\scriptsize [$-$0.12, 0.04]} & {\scriptsize [$-$0.23, 0.07]} & {\scriptsize [$-$0.29, 0.09]}\\
 &  & 6 & 639 & $-$0.08 & $-$0.05 & $-$0.03 & $-$0.03 & $-$0.12 & $-$0.21$^{*}$\\
 & & & & {\scriptsize [$-$0.21, 0.06]} & {\scriptsize [$-$0.20, 0.15]} & {\scriptsize [$-$0.18, 0.14]} & {\scriptsize [$-$0.11, 0.08]} & {\scriptsize [$-$0.28, 0.02]} & {\scriptsize [$-$0.43, $-$0.01]}\\
 &  & 12 & 633 & $-$0.22$^{*}$ & $-$0.11 & $-$0.06 & $-$0.05 & $-$0.13 & $-$0.21\\
 & & & & {\scriptsize [$-$0.39, $-$0.02]} & {\scriptsize [$-$0.27, 0.09]} & {\scriptsize [$-$0.24, 0.15]} & {\scriptsize [$-$0.16, 0.05]} & {\scriptsize [$-$0.32, 0.05]} & {\scriptsize [$-$0.46, 0.04]}\\
Real & real & 3 & 642 & 0.05 & 0.03 & 0.02 & $-$0.04 & $-$0.07 & $-$0.10\\
 & & & & {\scriptsize [$-$0.05, 0.19]} & {\scriptsize [$-$0.09, 0.17]} & {\scriptsize [$-$0.08, 0.12]} & {\scriptsize [$-$0.14, 0.05]} & {\scriptsize [$-$0.24, 0.07]} & {\scriptsize [$-$0.31, 0.10]}\\
 &  & 6 & 639 & 0.03 & 0.04 & 0.03 & $-$0.05 & $-$0.14 & $-$0.23$^{*}$\\
 & & & & {\scriptsize [$-$0.08, 0.14]} & {\scriptsize [$-$0.08, 0.17]} & {\scriptsize [$-$0.07, 0.14]} & {\scriptsize [$-$0.16, 0.05]} & {\scriptsize [$-$0.32, 0.02]} & {\scriptsize [$-$0.48, 0.00]}\\
 &  & 12 & 633 & 0.08 & 0.03 & 0.01 & $-$0.09 & $-$0.20 & $-$0.28$^{*}$\\
 & & & & {\scriptsize [$-$0.08, 0.27]} & {\scriptsize [$-$0.08, 0.15]} & {\scriptsize [$-$0.09, 0.12]} & {\scriptsize [$-$0.24, 0.03]} & {\scriptsize [$-$0.43, 0.01]} & {\scriptsize [$-$0.57, $-$0.01]}\\
\bottomrule\end{tabular}
\par\smallskip\parbox{0.96\linewidth}{\scriptsize As 2c. Per +1 pp on the 2-year yield. $^{*}$ = 90\% interval excludes zero; intervals in brackets. Obs = monthly observations in the equations.}
\end{table}


\section{Robustness 3: credit spreads}
\begin{table}[p]\centering\footnotesize\setlength{\tabcolsep}{2.2pt}
\caption{Robustness 3a: r*-gap FCI plus the Baa-10-year spread, start 1994}\label{tab:R3a}
\begin{tabular}{llcrrrrrrr}\toprule
 & & & & \multicolumn{3}{c}{Unemployment gap (pp)} & \multicolumn{3}{c}{Price level (\%)}\\ \cmidrule(lr){5-7}\cmidrule(lr){8-10}
2-year & FCI & Lags & Obs & 12 m & 24 m & 36 m & 12 m & 24 m & 36 m\\ \midrule
r*-gap & r*-gap + Baa spread & 3 & 392 & 0.23 & 0.14 & 0.00 & $-$0.14 & $-$0.16 & $-$0.15\\
 & & & & {\scriptsize [$-$0.04, 0.47]} & {\scriptsize [$-$0.10, 0.36]} & {\scriptsize [$-$0.17, 0.16]} & {\scriptsize [$-$0.32, 0.06]} & {\scriptsize [$-$0.53, 0.18]} & {\scriptsize [$-$0.63, 0.32]}\\
 &  & 6 & 392 & 0.26$^{*}$ & 0.26$^{*}$ & 0.11 & $-$0.22$^{*}$ & $-$0.44$^{*}$ & $-$0.62$^{*}$\\
 & & & & {\scriptsize [0.02, 0.54]} & {\scriptsize [0.01, 0.52]} & {\scriptsize [$-$0.07, 0.31]} & {\scriptsize [$-$0.40, $-$0.02]} & {\scriptsize [$-$0.85, $-$0.08]} & {\scriptsize [$-$1.25, $-$0.08]}\\
 &  & 12 & 392 & 0.21 & 0.38$^{*}$ & 0.23$^{*}$ & $-$0.28$^{*}$ & $-$0.61$^{*}$ & $-$0.84$^{*}$\\
 & & & & {\scriptsize [$-$0.20, 0.55]} & {\scriptsize [0.09, 0.64]} & {\scriptsize [0.00, 0.44]} & {\scriptsize [$-$0.47, $-$0.09]} & {\scriptsize [$-$0.99, $-$0.20]} & {\scriptsize [$-$1.49, $-$0.23]}\\
\bottomrule\end{tabular}
\par\smallskip\parbox{0.96\linewidth}{\scriptsize The spread enters the daily block and responds on impact (event-day change per +1 pp 2-year). Only with the FCI. Per +1 pp on the 2-year yield. $^{*}$ = 90\% interval excludes zero; intervals in brackets. Obs = monthly observations in the equations.}
\end{table}
\begin{table}[p]\centering\footnotesize\setlength{\tabcolsep}{2.2pt}
\caption{Robustness 3b: excess bond premium as a third monthly variable, start 1973}\label{tab:R3b}
\begin{tabular}{llcrrrrrrr}\toprule
 & & & & \multicolumn{3}{c}{Unemployment gap (pp)} & \multicolumn{3}{c}{Price level (\%)}\\ \cmidrule(lr){5-7}\cmidrule(lr){8-10}
2-year & FCI & Lags & Obs & 12 m & 24 m & 36 m & 12 m & 24 m & 36 m\\ \midrule
r*-gap & none & 3 & 642 & 0.01 & 0.04 & 0.04 & $-$0.09 & $-$0.15 & $-$0.19\\
 & & & & {\scriptsize [$-$0.10, 0.15]} & {\scriptsize [$-$0.10, 0.21]} & {\scriptsize [$-$0.09, 0.19]} & {\scriptsize [$-$0.25, 0.06]} & {\scriptsize [$-$0.52, 0.23]} & {\scriptsize [$-$0.72, 0.42]}\\
 &  & 6 & 639 & 0.01 & 0.03 & 0.01 & $-$0.12$^{*}$ & $-$0.31$^{*}$ & $-$0.51$^{*}$\\
 & & & & {\scriptsize [$-$0.10, 0.14]} & {\scriptsize [$-$0.14, 0.19]} & {\scriptsize [$-$0.19, 0.16]} & {\scriptsize [$-$0.26, $-$0.02]} & {\scriptsize [$-$0.61, $-$0.04]} & {\scriptsize [$-$1.06, $-$0.09]}\\
 &  & 12 & 633 & $-$0.03 & $-$0.01 & 0.01 & $-$0.23$^{*}$ & $-$0.50$^{*}$ & $-$0.78$^{*}$\\
 & & & & {\scriptsize [$-$0.17, 0.15]} & {\scriptsize [$-$0.19, 0.16]} & {\scriptsize [$-$0.19, 0.18]} & {\scriptsize [$-$0.38, $-$0.11]} & {\scriptsize [$-$0.83, $-$0.23]} & {\scriptsize [$-$1.32, $-$0.28]}\\
r*-gap & r*-gap & 3 & 642 & 0.03 & 0.02 & 0.00 & $-$0.09 & $-$0.19 & $-$0.28\\
 & & & & {\scriptsize [$-$0.09, 0.14]} & {\scriptsize [$-$0.13, 0.16]} & {\scriptsize [$-$0.16, 0.14]} & {\scriptsize [$-$0.22, 0.08]} & {\scriptsize [$-$0.51, 0.18]} & {\scriptsize [$-$0.75, 0.29]}\\
 &  & 6 & 639 & 0.04 & 0.03 & 0.01 & $-$0.12 & $-$0.31 & $-$0.55\\
 & & & & {\scriptsize [$-$0.09, 0.19]} & {\scriptsize [$-$0.15, 0.18]} & {\scriptsize [$-$0.18, 0.15]} & {\scriptsize [$-$0.26, 0.03]} & {\scriptsize [$-$0.62, 0.02]} & {\scriptsize [$-$1.04, 0.04]}\\
 &  & 12 & 633 & 0.07 & 0.04 & 0.04 & $-$0.24$^{*}$ & $-$0.56$^{*}$ & $-$0.86$^{*}$\\
 & & & & {\scriptsize [$-$0.11, 0.25]} & {\scriptsize [$-$0.17, 0.24]} & {\scriptsize [$-$0.16, 0.24]} & {\scriptsize [$-$0.40, $-$0.08]} & {\scriptsize [$-$0.89, $-$0.22]} & {\scriptsize [$-$1.37, $-$0.37]}\\
\bottomrule\end{tabular}
\par\smallskip\parbox{0.96\linewidth}{\scriptsize As 3a. Per +1 pp on the 2-year yield. $^{*}$ = 90\% interval excludes zero; intervals in brackets. Obs = monthly observations in the equations.}
\end{table}
\begin{table}[p]\centering\footnotesize\setlength{\tabcolsep}{2.2pt}
\caption{Robustness 3c: excess bond premium as a third monthly variable, start 1973}\label{tab:R3c}
\begin{tabular}{llcrrrrrrr}\toprule
 & & & & \multicolumn{3}{c}{Unemployment gap (pp)} & \multicolumn{3}{c}{Price level (\%)}\\ \cmidrule(lr){5-7}\cmidrule(lr){8-10}
2-year & FCI & Lags & Obs & 12 m & 24 m & 36 m & 12 m & 24 m & 36 m\\ \midrule
r*-gap & none & 3 & 642 & 0.01 & 0.04 & 0.04 & $-$0.09 & $-$0.15 & $-$0.19\\
 & & & & {\scriptsize [$-$0.10, 0.15]} & {\scriptsize [$-$0.10, 0.21]} & {\scriptsize [$-$0.09, 0.19]} & {\scriptsize [$-$0.25, 0.06]} & {\scriptsize [$-$0.52, 0.23]} & {\scriptsize [$-$0.72, 0.42]}\\
 &  & 6 & 639 & 0.01 & 0.03 & 0.01 & $-$0.12$^{*}$ & $-$0.31$^{*}$ & $-$0.51$^{*}$\\
 & & & & {\scriptsize [$-$0.10, 0.14]} & {\scriptsize [$-$0.14, 0.19]} & {\scriptsize [$-$0.19, 0.16]} & {\scriptsize [$-$0.26, $-$0.02]} & {\scriptsize [$-$0.61, $-$0.04]} & {\scriptsize [$-$1.06, $-$0.09]}\\
 &  & 12 & 633 & $-$0.03 & $-$0.01 & 0.01 & $-$0.23$^{*}$ & $-$0.50$^{*}$ & $-$0.78$^{*}$\\
 & & & & {\scriptsize [$-$0.17, 0.15]} & {\scriptsize [$-$0.19, 0.16]} & {\scriptsize [$-$0.19, 0.18]} & {\scriptsize [$-$0.38, $-$0.11]} & {\scriptsize [$-$0.83, $-$0.23]} & {\scriptsize [$-$1.32, $-$0.28]}\\
r*-gap & r*-gap & 3 & 642 & 0.03 & 0.02 & 0.00 & $-$0.09 & $-$0.19 & $-$0.28\\
 & & & & {\scriptsize [$-$0.09, 0.14]} & {\scriptsize [$-$0.13, 0.16]} & {\scriptsize [$-$0.16, 0.14]} & {\scriptsize [$-$0.22, 0.08]} & {\scriptsize [$-$0.51, 0.18]} & {\scriptsize [$-$0.75, 0.29]}\\
 &  & 6 & 639 & 0.04 & 0.03 & 0.01 & $-$0.12 & $-$0.31 & $-$0.55\\
 & & & & {\scriptsize [$-$0.09, 0.19]} & {\scriptsize [$-$0.15, 0.18]} & {\scriptsize [$-$0.18, 0.15]} & {\scriptsize [$-$0.26, 0.03]} & {\scriptsize [$-$0.62, 0.02]} & {\scriptsize [$-$1.04, 0.04]}\\
 &  & 12 & 633 & 0.07 & 0.04 & 0.04 & $-$0.24$^{*}$ & $-$0.56$^{*}$ & $-$0.86$^{*}$\\
 & & & & {\scriptsize [$-$0.11, 0.25]} & {\scriptsize [$-$0.17, 0.24]} & {\scriptsize [$-$0.16, 0.24]} & {\scriptsize [$-$0.40, $-$0.08]} & {\scriptsize [$-$0.89, $-$0.22]} & {\scriptsize [$-$1.37, $-$0.37]}\\
\bottomrule\end{tabular}
\par\smallskip\parbox{0.96\linewidth}{\scriptsize As 3b. Per +1 pp on the 2-year yield. $^{*}$ = 90\% interval excludes zero; intervals in brackets. Obs = monthly observations in the equations.}
\end{table}


\end{document}
